% ============================================================
% CAPÍTULO 3 — Estado Actual del Sistema
% ============================================================
\chapter{Estado Actual del Sistema}
\label{cap:estado}

\section{Componentes funcionales verificados}

A fecha de redacci\'on del presente informe, los siguientes componentes
del sistema operan de forma verificada en entorno de pruebas:

\begin{itemize}[leftmargin=*]
    \item \textbf{Telemetr\'ia extremo a extremo.} Un nodo sensor
    ESP32-S3 despierta de \textit{deep sleep}, lee sus sondas DS18B20 y
    sus sensores capacitivos de humedad de suelo, emite un
    \texttt{SENSOR\_CLUSTER\_REPORT} por ESP-NOW al \textit{gateway},
    \'este lo reenv\'ia por UART a la Raspberry Pi, que lo decodifica y lo
    publica por WebSocket al backend, donde queda persistido en
    PostgreSQL y visible en el dashboard.
    \item \textbf{Comando manual extremo a extremo.} La cadena
    dashboard $\to$ \texttt{POST /api/command} $\to$ Redis Pub/Sub $\to$
    \textit{dispatcher} WebSocket $\to$ Raspberry Pi $\to$ UART $\to$
    \textit{gateway} $\to$ ESP-NOW $\to$ nodo actuador opera de forma
    completa para el comando \texttt{SET\_GPIO}.
    \item \textbf{Handshake y enrutado.} Establecimiento de sesi\'on
    SYN / SYN-ACK / ACK entre nodo y \textit{gateway}, y alta de rutas en
    caliente mediante \texttt{ROUTE\_ADD}.
    \item \textbf{Secuencias programadas.} El servicio secuenciador
    ejecuta secuencias multi-paso definidas desde el planificador del
    frontend.
    \item \textbf{Trazabilidad LIMS.} Alta de plantas y experimentos,
    asignaci\'on N:M entre ambos y consulta de telemetr\'ia filtrada por
    planta.
    \item \textbf{Autenticaci\'on.} Registro e inicio de sesi\'on con JWT
    y \texttt{bcrypt}, con protecci\'on efectiva de las rutas privadas.
    \item \textbf{SDK cient\'ifico.} Extracci\'on de series a
    \texttt{DataFrame} con cach\'e Parquet, limpieza y remuestreo,
    c\'alculos agron\'omicos vectorizados y exportaci\'on a CSV y Excel.
    \item \textbf{Despliegue reproducible.} Los cinco entornos de
    \textit{Docker Compose} (desarrollo, \textit{override},
    \textit{staging}, producci\'on y Raspberry Pi) levantan desde un
    \'arbol limpio, y la integraci\'on continua construye y publica las
    im\'agenes de backend y frontend en GitHub Container Registry.
\end{itemize}

\textbf{Escenario de prueba m\'as representativo}. La validaci\'on
integrada se realiz\'o sobre banco de laboratorio con \textbf{tres nodos
ESP32-S3} (uno \textit{gateway} conectado por UART a la Raspberry Pi y
dos nodos de campo) y hasta dos plantas instrumentadas por nodo, con un
ciclo de \textit{deep sleep} de 60 segundos. El procedimiento completo,
incluido el conexionado, est\'a documentado en
\texttt{docs/deep\_sleep\_test.md} y es reproducible. Es importante
precisar el alcance de esta validaci\'on: se trata de \textbf{sesiones de
prueba en laboratorio, no de una campa\~na de operaci\'on continua}. El
sistema nunca lleg\'o a desplegarse en un invernadero real ni a
funcionar sin supervisi\'on durante d\'ias.

\section{Componentes parciales o no implementados}

Por honestidad t\'ecnica y para facilitar la replicabilidad o la
continuaci\'on parcial del proyecto, se documentan los aspectos que se
encuentran en estado parcial o pendiente:

\begin{itemize}[leftmargin=*]
    \item \textbf{Cifrado en la capa de campo.} AES-128 sobre el
    \textit{payload} ESP-NOW estaba planificado para la versi\'on 3.0 y
    no est\'a implementado. La trama es legible por cualquier ESP32
    cercano sintonizado al canal.
    \item \textbf{Discovery din\'amico incompleto.} Existe un
    \textit{Device Manager} con cach\'e persistente de topolog\'ia y un
    \textit{router} de descubrimiento en el backend, pero el alta de un
    nodo nuevo no es aut\'onoma: la direcci\'on MAC del \textit{gateway}
    est\'a fijada en tiempo de compilaci\'on en el firmware de los nodos y
    el \texttt{NODE\_ID} se asigna tambi\'en por \textit{build flag}.
    A\~nadir un nodo exige recompilar y reflashear.
    \item \textbf{Sin observabilidad.} No hay agregaci\'on de
    \textit{logs}, m\'etricas, alarmas ni notificaciones. El diagn\'ostico
    de una incidencia obliga a leer la salida de cada contenedor por
    separado.
    \item \textbf{Dashboard no adaptado a m\'ovil.} La interfaz est\'a
    dise\~nada para pantalla de escritorio.
    \item \textbf{Riego f\'isico sin validar.} La instalaci\'on
    hidr\'aulica adquirida (cap\'itulo~\ref{cap:sostenibilidad}) no lleg\'o
    a montarse. El \textit{spike} de validaci\'on de bombeo
    (\texttt{docs/spikes/SPIKE-001}) qued\'o escrito y sin ejecutar, y
    adem\'as est\'a redactado para una bomba de 5\,V mientras que el
    material comprado es de 12\,V, por lo que su plan de pruebas
    tampoco ser\'ia aplicable sin reescribirlo.
\end{itemize}

Conviene se\~nalar que \textbf{el c\'odigo no contiene marcas
\texttt{TODO} ni \texttt{FIXME} significativas}: no hay funciones a
medio escribir ni parches provisionales pendientes de resolver. Lo que
falta no est\'a dentro de lo construido, sino que son capas que nunca se
empezaron. Es deuda de alcance, no de implementaci\'on.

Se considera que ninguna de estas limitaciones impide la utilidad
docente ni el aprovechamiento del sistema como caso de estudio: la
arquitectura est\'a documentada y los puntos parciales est\'an claramente
identificados.

\section{Tests y validaci\'on realizados}

La validaci\'on del sistema se ha articulado en m\'ultiples niveles.

\subsection{Tests unitarios por capa}

Cada capa cuenta con suite de tests unitarios independiente:

\begin{table}[H]
    \centering
    \caption{Tests unitarios por capa.}
    \label{tab:cobertura_tests}
    \begin{tabular}{llc}
        \toprule
        \textbf{Capa} & \textbf{Framework} & \textbf{N\'umero de tests} \\
        \midrule
        Firmware ESP32        & Unity (C++)     & 32 \\
        Servicio Raspberry Pi & pytest          & 37 \\
        Backend FastAPI       & pytest          & 16 \\
        Frontend              & Vitest          & 3 \\
        SDK cient\'ifico      & pytest          & 52 \\
        \midrule
        \textbf{TOTAL}        &                 & \textbf{140} \\
        \bottomrule
    \end{tabular}
\end{table}

\textbf{Sobre la cobertura}. La tabla no incluye porcentaje de cobertura
porque \textbf{el proyecto nunca instrument\'o su medida}: no hay
configuraci\'on de \texttt{coverage} ni umbral m\'inimo en integraci\'on
continua, y ning\'un informe hist\'orico que citar. Se ha preferido
declararlo abiertamente antes que producir ahora una cifra retrospectiva
que no represente el estado del proyecto durante su desarrollo. El
reparto de los tests es, en s\'i mismo, informativo: las capas con
l\'ogica algor\'itmica pura y f\'acilmente aislable (SDK cient\'ifico,
protocolo del firmware, cach\'e de telemetr\'ia) concentran la mayor
parte de las pruebas, mientras que las capas dominadas por integraci\'on
con servicios externos (backend, frontend) quedan claramente por debajo.

\subsection{Tests de integraci\'on}

Adem\'as de los tests unitarios se mantienen:

\begin{itemize}[leftmargin=*]
    \item Un entorno de integraci\'on del firmware
    (\texttt{pio test -e integration}) que ejercita los flujos completos
    del \texttt{SystemManager} y del \texttt{ProtocolEngine} sobre
    \textit{mocks} de Arduino, sin hardware.
    \item Tests de backend sobre cliente HTTP de prueba que recorren
    autenticaci\'on, emisi\'on de comandos, LIMS y an\'alisis contra una
    base de datos real levantada en contenedor.
    \item Validaci\'on manual extremo a extremo sobre el entorno de
    \textit{staging}, con la Raspberry Pi conectada al backend remoto y
    los tres nodos operativos.
\end{itemize}

\subsection{Pruebas de campo}

No se realizaron. El sistema se valid\'o \'unicamente sobre banco de
laboratorio, en el escenario descrito al principio de este cap\'itulo.
No hubo instalaci\'on en invernadero, ni operaci\'on desatendida
prolongada, ni exposici\'on de los nodos a condiciones ambientales
reales.

\subsection{M\'etricas medidas}

\textbf{El sistema no lleg\'o a caracterizarse cuantitativamente.} No se
midieron la latencia extremo a extremo, la tasa de p\'erdida de tramas
ESP-NOW, la vida \'util real de un nodo alimentado por bater\'ia, el
\textit{throughput} sostenido de telemetr\'ia ni el tiempo medio de
reconexi\'on del WebSocket. Los \'unicos valores de consumo que aparecen
en este informe ($\sim$80\,mA en activo, $\sim$0,01\,mA en reposo)
proceden de la hoja de caracter\'isticas del ESP32-S3, no de una medida
propia.

Se ha optado por declarar esta ausencia en lugar de rellenar una tabla
con estimaciones. La falta de caracterizaci\'on no es un descuido
aislado: es exactamente el tipo de trabajo que un proyecto sostenido por
una sola persona en horario nocturno termina posponiendo, y como tal
refuerza el an\'alisis del cap\'itulo~\ref{cap:sostenibilidad} en lugar de
contradecirlo. Cualquier continuaci\'on del proyecto deber\'ia empezar
por aqu\'i: sin l\'inea base medida no hay forma de saber si un cambio
mejora o empeora el sistema.
